\documentclass[11pt,a4paper]{article}
\usepackage[T1]{fontenc}
\usepackage[utf8]{inputenc}
\usepackage{lmodern}
\usepackage{amsmath,amssymb,amsthm,mathtools}
\usepackage[margin=1.05in]{geometry}
\usepackage{xcolor}
\usepackage{hyperref}
\usepackage{enumitem}
\usepackage{booktabs}
\hypersetup{colorlinks=true,linkcolor=blue!50!black,citecolor=blue!50!black,urlcolor=blue!50!black}

\newtheorem{theorem}{Theorem}[section]
\newtheorem{lemma}[theorem]{Lemma}
\newtheorem{proposition}[theorem]{Proposition}
\newtheorem{corollary}[theorem]{Corollary}
\theoremstyle{definition}
\newtheorem{definition}[theorem]{Definition}
\theoremstyle{remark}
\newtheorem{remark}[theorem]{Remark}
\newtheorem{convention}[theorem]{Convention}

\DeclareMathOperator{\Tr}{Tr}
\DeclareMathOperator{\Cov}{Cov}
\DeclareMathOperator{\artanh}{artanh}
\newcommand{\R}{\mathbb{R}}
\newcommand{\C}{\mathbb{C}}
\newcommand{\Z}{\mathbb{Z}}
\newcommand{\N}{\mathbb{N}}
\newcommand{\cA}{\mathcal{A}}
\newcommand{\cP}{\mathcal{P}}
\newcommand{\cX}{\mathcal{X}}
\newcommand{\dU}{\,dU}
\newcommand{\norm}[1]{\lVert #1\rVert}
\newcommand{\Fref}[3]{\cite[#1~`#2' (\texttt{#3})]{F}}
\newcommand{\Flab}[2]{\cite[#1 \texttt{#2}]{F}}
\newcommand{\Fname}[2]{\cite[#1~`#2']{F}}
\newcommand{\IIref}[3]{\cite[#1~`#2' (\texttt{#3})]{II}}
\newcommand{\IIlab}[2]{\cite[#1 \texttt{#2}]{II}}
\newcommand{\IIIref}[3]{\cite[#1~`#2' (\texttt{#3})]{III}}
\newcommand{\IIIlab}[2]{\cite[#1 \texttt{#2}]{III}}

\title{Lattice Yang--Mills theory at strong coupling IV:\\ a mass-gap bound with the correct leading coefficient\\
\large A sequel to the foundations monograph of the project ``Solving Yang--Mills and the Mass Gap''}
\author{}
\date{9 November 2026; revised the same day after an independent referee report; living document}

\begin{document}
\maketitle

\begin{abstract}
For lattice Yang--Mills theory with Wilson's action in dimension $d\ge3$ and $0<\beta<\beta_0(d)$ (the threshold of
the cluster expansion of [F]), we prove that the mass gap of the infinite-volume theory of Part~II, and the finite-volume
gaps of [F] on spatial lattices with all sides $\ge4$, are at least $4\log(1/q(\beta))$, with $q(\beta)=e^3K(e^{4\beta}-1)$ and
$K=4(2d-3)$. As $\beta\downarrow0$ this is $4\log(1/\beta)-O(1)$. The leading coefficient $4$ is the one predicted by the
strong-coupling expansion (the lightest glueball is a tube with four plaquettes per unit of time, with predicted mass
$-4\log u(\beta)+o(1)$, $u(\beta)\sim\beta/4$ for $SU(2)$), whereas the Dobrushin bounds of [F] and Part~II give only
$\log(1/\beta)-O(1)$. The
proof combines three ingredients: duplicate variables, by which every term of the expansion of a covariance that does
not connect the two observables cancels by the swap symmetry; the loop expansion and the ratio bound of Part~III for the
duplicated system; and a geometric lemma: a surface without free edges has at least four temporal plaquettes in each
time layer it crosses, because these plaquettes form a graph of minimum degree two on the spatial lattice, whose girth
is four (sides $\ge4$ are needed: with a side of length $3$ the girth is three). In $d=2$ the girth is the spatial size,
which gives an infinite mass gap in infinite volume and a finite-volume gap growing linearly with the spatial size, in
agreement with the exact solution. As in [F], II and III nothing non-abelian is used, so the result belongs to the
strong-coupling endpoint of the problem; it is the lower half of what principle~3 (``the strong-coupling test'') of the
research program requires, the matching upper bound being open.
\end{abstract}

\tableofcontents

\section{Introduction and main result}\label{mg:sec:intro}

The foundations monograph \cite{F} proves that the finite-volume mass gap is bounded below uniformly in the volume for
small $\beta$, by Dobrushin's method; Part~II \cite{II} constructs the infinite-volume theory, identifies its transfer
operator with a reversible Markov chain of time slices, and proves the same kind of bound,
$m\ge\log(1/(6(d-1)\tanh\beta))$, for the infinite-volume mass gap $m(\mu_\beta)$. These bounds behave like
$\log(1/\beta)$ as $\beta\downarrow0$. The strong-coupling expansion predicts $m(\beta)=-4\log u(\beta)+O(u)$, where
$u(\beta)=c_\rho(\beta)/(Nc_0(\beta))$ ($\approx\beta/4$ for $SU(2)$, $\approx\beta/(2N^2)$ for $N\ge3$): the lightest
state in the gauge-invariant sector propagates as a tube of plaquettes with four plaquettes per unit of Euclidean time.
Principle~3 of the research program (\texttt{docs/03-research-program.md}, \S3.1) asks every argument in the
strong-coupling regime to reproduce this behaviour. Here we prove a lower bound with the correct leading coefficient.

\emph{Literature.} The strong-coupling glueball spectrum was studied rigorously in the 1980s: by their abstracts, Schor
\cite{Sch83} proved the existence of glueballs in strongly coupled lattice gauge theories, and O'Carroll, Braga and Schor
\cite{OBS85} (in $2+1$ dimensions) and O'Carroll and Barbosa \cite{OB85} (in $3+1$ dimensions) obtained isolated masses of the
form $-4\log\beta+r(\beta)$. We have not yet read these papers (they are marked ``verify''), and nothing below relies on them.
What this article adds for the project is a short self-contained proof with explicit constants and range, for the
infinite-volume theory of \cite{II}, and a covariance bound for all bounded local observables.

\begin{convention}[Labels]
As in \Fname{Convention}{Labels}: every Theorem, Proposition, Lemma and Corollary below has a complete proof here or in
\cite{F,II,III}; Remarks are commentary. We use the notation of \cite{III}: $K:=4(2d-3)$,
$\epsilon(\beta)=e^{2\beta}-1$, $x(\beta)=e^2K\epsilon(\beta)$, $\theta(\beta)=x/(1-x)$, and
$\beta_0(d)=\frac12\log(1+\frac1{e^2K(K+2)})$, the threshold below which \Fref{Theorem}{Convergence of the
strong-coupling cluster expansion}{thm:conv} holds.
\end{convention}

\begin{theorem}[Main theorem]\label{mg:thm:main}
Let $G,\rho$ be as in \Fname{Convention}{Group} with $G$ metrisable, $d\ge3$, $0<\beta<\beta_0(d)$, and
\[
q(\beta):=e^3K\bigl(e^{4\beta}-1\bigr).
\]
Then $q(\beta)<\frac{e(2+\epsilon_0)}{K+2}<0.389$, where $\epsilon_0=\frac1{e^2K(K+2)}$, and:
\begin{enumerate}[label=(\roman*)]
\item for every periodic spatial lattice $\Lambda_s$ with all sides $L_i\ge4$, the finite-volume mass gap of
\Fref{Definition}{Hamiltonian and mass gap}{def:H} satisfies $m_{\Lambda_s}(\beta)\ge4\log(1/q(\beta))$;
\item the mass gap of the infinite-volume theory \IIref{Definition}{Mass gap of an infinite-volume lattice theory}{iv:def:gap}
satisfies $m(\mu_\beta)\ge4\log(1/q(\beta))$;
\item for bounded measurable $A,B$ on $G^{E(\Z^d)}$ depending only on finite link sets $S_A,S_B$, all of whose
endpoints have time coordinates in $[a_1,a_2]$, resp.\ $[b_1,b_2]$, with $a_2<b_1$,
\[
\bigl|\Cov_{\mu_\beta}(A,B)\bigr|\le\frac{2\norm A_\infty\norm B_\infty}{1-q(\beta)}\,
c(\beta)^{|S_A|+|S_B|}\;q(\beta)^{4(b_1-a_2)},
\]
where $c(\beta):=\bigl(e^{2\theta(\beta)}(1+\frac1K)\bigr)^{2(d-1)}$.
\end{enumerate}
As $\beta\downarrow0$, $4\log(1/q(\beta))=4\log(1/\beta)-4\log(4e^3K)+O(\beta)$.
\end{theorem}

Here $\Cov(A,B):=\mu(AB)-\mu(A)\mu(B)$. Statement (iii) is the infinite-volume form of the covariance bound of
Theorem~\ref{mg:thm:finite}; (i) follows from Theorem~\ref{mg:thm:finite} by \Fref{Corollary}{Gap from Euclidean
correlations}{cor:gapfromcorr}, and (ii) from (iii) by \IIref{Theorem}{OS reconstruction of the lattice theory}{iv:thm:OS}(iii).

\begin{remark}[Comparison]\label{mg:rem:comparison}
For $d=4$: the bound of Part~II, $\log(1/(18\tanh\beta))$, is $\log(1/\beta)-2.89+O(\beta^2)$; the bound here is
$4\log(1/\beta)-29.5+O(\beta)$. In the common range of validity $\beta<\beta_0(4)\approx1.54\cdot10^{-4}$ the two are
comparable near $\beta_0(4)$ ($5.890$ and $5.591$ at $\beta_0(4)$), and the present one is larger for
$\beta<1.39\cdot10^{-4}$, by an additive amount that tends to infinity as $\beta\to0$ (the ratio of the two bounds increases
to $4$; it is $1.92$ at $\beta=10^{-5}$). These values were computed with \texttt{python3}. The point of the theorem is the
coefficient $4$, not the numerical strength at fixed $\beta$; the additive constant $-29.5$ is far from the predicted
$+4\log4$ for $SU(2)$ (from $u\approx\beta/4$).
\end{remark}

\paragraph{Strategy.} (1) For the product measure $\mu\otimes\mu$ of two independent copies $(U,U')$,
$2\Cov_\mu(A,B)=\int(A-A')(B-B')\,d(\mu\otimes\mu)$. The duplicated Wilson weight is
$\prod_p(1+f_\beta(U_p))(1+f_\beta(U'_p))=\prod_p(1+g_p)$ with $|g_p|\le e^{4\beta}-1$, so the loop expansion of
\cite{III} applies to the duplicated system, with partition functions $\Xi(\cdot)^2$. (2) In this expansion every set of
plaquettes that does not contain a component connecting the supports of $A$ and $B$ gives a product of two integrals, each
odd under the exchange $U\leftrightarrow U'$, hence zero (Lemma~\ref{mg:lem:cancel}). (3) A connecting component has no
free edges away from the supports, so in every time layer between the supports it contains at least four temporal plaquettes
(Lemma~\ref{mg:lem:layer}), hence at least $4\times$(number of layers) plaquettes (Corollary~\ref{mg:cor:size}). (4) The sum
over such sets is controlled by the ratio bound and the rooted counting of \cite{III}.

\section{The duplicated system}\label{mg:sec:dup}

Throughout this section and the next, $\Lambda=\prod_{\lambda=0}^{d-1}\Z/L_\lambda\Z$ is a periodic lattice with all
$L_\lambda\ge3$, with the conventions of \cite{F,III}: $\cP=\cP(\Lambda)$, $E(\Lambda)$, $E(p)$, $E(X)$, adjacency and
connectedness of sets of plaquettes, $f_\beta=w_\beta/c_0(\beta)-1$, and for $\mathcal Q\subseteq\cP$ the restricted
partition function $\Xi(\mathcal Q)=\int\prod_{p\in\mathcal Q}(1+f_\beta(U_p))\dU$ of \IIIref{Definition}{Restricted partition
functions, loop polymers}{al:def:Xi}, so that $\Xi(\cP)=Z_\Lambda(\beta)/c_0(\beta)^{|\cP|}$.

\begin{definition}[Duplicated system]\label{mg:def:dup}
On $\Omega_\Lambda\times\Omega_\Lambda$ with the product Haar measure $dU\,dU'$ put, for $p\in\cP$,
\[
g_p(U,U'):=\bigl(1+f_\beta(U_p)\bigr)\bigl(1+f_\beta(U'_p)\bigr)-1 .
\]
For $X\subseteq\cP$ and bounded measurable $F$ on $\Omega_\Lambda\times\Omega_\Lambda$ let
$\Phi_F(X):=\iint F\prod_{p\in X}g_p\,dU\,dU'$.
\end{definition}

\begin{lemma}[Properties of the duplicated system]\label{mg:lem:dup}
\begin{enumerate}[label=(\alph*)]
\item $\norm{g_p}_\infty\le\epsilon(2\beta)=e^{4\beta}-1$.
\item For $\mathcal Q\subseteq\cP$, $\iint\prod_{p\in\mathcal Q}(1+g_p)\,dU\,dU'=\Xi(\mathcal Q)^2$; and
$\mu_{\Lambda,\beta}\otimes\mu_{\Lambda,\beta}$ has density $\prod_{p\in\cP}(1+g_p)/\Xi(\cP)^2$ with respect to $dU\,dU'$.
\item Let $X\subseteq\cP$, let $\ell$ be a link lying in exactly one plaquette of $X$, and let $H$ be bounded measurable and
independent of $(U_\ell,U'_\ell)$. Then $\iint H\prod_{p\in X}g_p\,dU\,dU'=0$.
\end{enumerate}
\end{lemma}

\begin{proof}
(a) By \IIIlab{Lemma}{al:lem:facts}(a),(b), $e^{-\beta}\le w_\beta\le e^\beta$ and $1\le c_0\le e^\beta$, so
$1+f_\beta=w_\beta/c_0\in[e^{-2\beta},e^{\beta}]$ and
$(1+f_\beta(U_p))(1+f_\beta(U'_p))\in[e^{-4\beta},e^{2\beta}]$ and $g_p\in[e^{-4\beta}-1,e^{2\beta}-1]$; both
$1-e^{-4\beta}$ and $e^{2\beta}-1$ are at most $e^{4\beta}-1$.
(b) The integrand is $\prod_{p\in\mathcal Q}(1+f_\beta(U_p))\cdot\prod_{p\in\mathcal Q}(1+f_\beta(U'_p))$; integrate over $U$ and $U'$
separately. For the density: $W_\beta(U)W_\beta(U')=c_0^{2|\cP|}\prod_p(1+g_p)$ and $Z_\Lambda^2=c_0^{2|\cP|}\Xi(\cP)^2$.
(c) Let $p$ be the plaquette of $X$ containing $\ell$. Then $U_p=AU_\ell^{\pm1}B$ and $U'_p=A'U_\ell'^{\pm1}B'$ with $A,B,A',B'$
independent of $(U_\ell,U'_\ell)$, and no other factor of the integrand depends on $(U_\ell,U'_\ell)$. By Fubini it suffices that
$\iint g_p\,dU_\ell\,dU'_\ell=\int(1+f_\beta(AU^{\pm1}B))\dU\cdot\int(1+f_\beta(A'U^{\pm1}B'))\dU-1=1\cdot1-1=0$, by
\IIIlab{Lemma}{al:lem:facts}(c).
\end{proof}

\begin{definition}[Loop polymers for a set of links]\label{mg:def:loop}
For $S\subseteq E(\Lambda)$ let $\cP_S:=\{p:E(p)\cap S\neq\emptyset\}$, let $\cX_S$ be the family of sets $X\subseteq\cP$ each of
whose connected components contains a plaquette of $\cP_S$ (including $X=\emptyset$), and for $X\subseteq\cP$ let
$\overline X:=\{p:E(p)\cap(E(X)\cup S)\neq\emptyset\}$. Put $J_S:=2(d-1)|S|$. Then $|\cP_S|\le J_S$ and
$|\overline X|\le(K+1)|X|+J_S$ (each link lies in $2(d-1)$ plaquettes, and each plaquette is adjacent to at most $K$ plaquettes,
\IIIlab{Lemma}{al:lem:facts}(d)).
\end{definition}

\begin{lemma}[Duplicated loop expansion]\label{mg:lem:loopexp}
Let $S\subseteq E(\Lambda)$ and let $F$ be bounded, measurable, and dependent only on the variables $(U_\ell,U'_\ell)$,
$\ell\in S$. Then
\[
\int F\,d(\mu_{\Lambda,\beta}\otimes\mu_{\Lambda,\beta})=\sum_{X\in\cX_S}\Phi_F(X)\,\Bigl(\frac{\Xi(\cP\setminus\overline X)}{\Xi(\cP)}\Bigr)^2 .
\]
\end{lemma}

\begin{proof}
This is the proof of \IIIref{Lemma}{Loop expansion}{al:lem:loopexp} with $f_\beta$ replaced by $g$, single link variables by
pairs $(U_\ell,U'_\ell)$, $dU$ by $dU\,dU'$, and $E_C$ by $S$. By Lemma~\ref{mg:lem:dup}(b),
$\Xi(\cP)^2\int F\,d(\mu\otimes\mu)=\sum_{T\subseteq\cP}\iint F\prod_{p\in T}g_p\,dU\,dU'$. For $T\subseteq\cP$ let $X(T)$ be the union of
the components of $T$ meeting $\cP_S$ and $Y(T):=T\setminus X(T)$. The proof of the claim in \IIIlab{Lemma}{al:lem:loopexp}, which
uses only the combinatorics of components, of $\cP_S$ (there $\cP_C$) and of $\overline X$, shows that $T\mapsto(X(T),Y(T))$ is a
bijection onto the pairs with $X\in\cX_S$ and $Y\subseteq\cP\setminus\overline X$. For such a pair, $F\prod_{p\in X}g_p$ depends only on
the variables of the links in $S\cup E(X)$ and $\prod_{p\in Y}g_p$ only on those in $E(Y)$, a disjoint set; so the integral factorises,
and $\sum_{Y\subseteq\cP\setminus\overline X}\iint\prod_{p\in Y}g_p=\iint\prod_{p\in\cP\setminus\overline X}(1+g_p)=\Xi(\cP\setminus\overline X)^2$
by Lemma~\ref{mg:lem:dup}(b). Divide by $\Xi(\cP)^2$.
\end{proof}

\begin{lemma}[Cancellation]\label{mg:lem:cancel}
Let $A,B$ be bounded measurable functions on $\Omega_\Lambda$ depending only on disjoint link sets $S_A,S_B$, put $S:=S_A\cup S_B$ and
$F(U,U'):=(A(U)-A(U'))(B(U)-B(U'))$. Then:
\begin{enumerate}[label=(\roman*)]
\item $\int F\,d(\mu\otimes\mu)=2\Cov_\mu(A,B)$ for every probability measure $\mu$ on $\Omega_\Lambda$, and $\norm F_\infty\le4\norm A_\infty\norm B_\infty$;
\item if $X\in\cX_S$ and no connected component of $X$ contains both a plaquette of $\cP_{S_A}$ and a plaquette of $\cP_{S_B}$, then
$\Phi_F(X)=0$;
\item if some link of $E(X)\setminus S$ lies in exactly one plaquette of $X$, then $\Phi_F(X)=0$.
\end{enumerate}
\end{lemma}

\begin{proof}
(i) Expand: $\int(A-A')(B-B')\,d(\mu\otimes\mu)=2\mu(AB)-2\mu(A)\mu(B)$.
(ii) Let $X_1$ be the union of the components of $X$ that contain a plaquette of $\cP_{S_A}$, and $X_2:=X\setminus X_1$; by
hypothesis no component of $X_1$ meets $\cP_{S_B}$, and every component of $X_2$ meets $\cP_{S_B}$ but not $\cP_{S_A}$. A component
containing a link of $S_B$ contains a plaquette of $\cP_{S_B}$; hence $E(X_1)\cap S_B=\emptyset$, and likewise
$E(X_2)\cap S_A=\emptyset$; and $E(X_1)\cap E(X_2)=\emptyset$, since distinct components share no link. So $S_A\cup E(X_1)$ and
$S_B\cup E(X_2)$ are disjoint, and since $dU\,dU'$ is a product measure,
\[
\Phi_F(X)=\iint(A-A')\prod_{p\in X_1}g_p\;\cdot\;\iint(B-B')\prod_{p\in X_2}g_p .
\]
The exchange $(U,U')\mapsto(U',U)$ preserves $dU\,dU'$ and every $g_p$ (which is symmetric in $U_p,U'_p$), and changes the sign of
$A-A'$; so the first factor equals its negative and vanishes.
(iii) Apply Lemma~\ref{mg:lem:dup}(c) with $H=F\prod_{p'\in X\setminus\{p\}}g_{p'}$, $p$ the plaquette containing the link.
\end{proof}

\section{Geometry: four plaquettes per time layer}\label{mg:sec:geometry}

Time coordinates are taken in $\Z/L_0\Z$. A link $(x,0)$ is \emph{temporal of layer} $x_0$ (it joins times $x_0$ and $x_0+1$); a
plaquette $(x,0i)$, $1\le i\le d-1$, is \emph{temporal of layer} $x_0$. For $Y\subseteq\cP$ and a layer $s$, $Y_s$ denotes the set of
temporal plaquettes of layer $s$ in $Y$, and $T(p)$ the set of time coordinates of the corners of $p$ (one time for a spatial
plaquette, two consecutive times for a temporal one). Let $\Lambda_s:=\prod_{i=1}^{d-1}\Z/L_i\Z$ be the spatial torus and $g$ the
\emph{girth} of its nearest-neighbour graph (the length of a shortest cycle).

\begin{lemma}[Girth]\label{mg:lem:girth}
Let all $L_i\ge3$. If $d\ge3$, then $g=4$ when all $L_i\ge4$, and $g=3$ when some $L_i=3$. If $d=2$, then $g=L_1$.
\end{lemma}

\begin{proof}
Since all $L_i\ge3$ the graph is simple. A cycle is a closed walk without backtracking, i.e.\ a sum of unit vectors $\pm e_i$ equal to
$0$ in $\Lambda_s$ with no two consecutive steps opposite. For $d=2$ the graph is the cycle $\Z/L_1\Z$ itself. For $d\ge3$ a unit square
is a cycle of length $4$; a cycle of length $3$ needs $\pm e_i\pm e_j\pm e_k=0$ with no two opposite consecutive steps, which forces
$i=j=k$ and $3e_i=0$, i.e.\ $L_i=3$; and conversely $L_i=3$ gives the cycle $0,e_i,2e_i$.
\end{proof}

\begin{lemma}[Plaquettes per layer]\label{mg:lem:layer}
Let all $L_i\ge3$ ($1\le i\le d-1$). Let $S\subseteq E(\Lambda)$ and $X\subseteq\cP$ be such that every link of $E(X)\setminus S$ lies in
at least two plaquettes of $X$. Let $s$ be a layer none of whose temporal links belongs to $S$, and let $Y$ be a connected component
of $X$ with $Y_s\neq\emptyset$. Then $|Y_s|\ge g$. In particular $|Y_s|\ge4$ if $d\ge3$ and all $L_i\ge4$.
\end{lemma}

\begin{proof}
The temporal links of layer $s$ are $\ell_{\vec x}:=((s,\vec x),0)$, $\vec x\in\Lambda_s$, and the plaquettes containing $\ell_{\vec x}$
are the temporal plaquettes $((s,\vec x),0i)$ and $((s,\vec x-e_i),0i)$ of layer $s$. Let $\Gamma$ be the graph on $\Lambda_s$ whose edges
are the plaquettes $((s,\vec x),0i)\in Y_s$, each joining $\vec x$ and $\vec x+e_i$ (the positions of its two temporal links). Since
$L_i\ge3$, distinct plaquettes give distinct edges and there are no loops, so $\Gamma$ is a subgraph of the nearest-neighbour graph of
$\Lambda_s$. If $\vec x$ is an endpoint of an edge, then $\ell_{\vec x}\in E(Y)\setminus S$, so $\ell_{\vec x}$ lies in at least two
plaquettes of $X$; these are temporal plaquettes of layer $s$ sharing the link $\ell_{\vec x}$ with a plaquette of $Y$, hence they belong
to the component $Y$, i.e.\ to $Y_s$. So every non-isolated vertex of $\Gamma$ has degree at least $2$, and since $Y_s\neq\emptyset$,
$\Gamma$ contains a cycle (follow edges without immediately backtracking until a vertex repeats), which has at least $g$ edges; they are
distinct plaquettes of $Y_s$. The last statement is Lemma~\ref{mg:lem:girth}.
\end{proof}

\begin{lemma}[Crossing layers]\label{mg:lem:crossing}
Let $Y\subseteq\cP$ be connected and let $s\neq s'$ be layers with $Y_s=Y_{s'}=\emptyset$. Then either $T(p)\subseteq I_1$ for all
$p\in Y$, or $T(p)\subseteq I_2$ for all $p\in Y$, where $I_1:=\{s+1,\dots,s'\}$ and $I_2:=\{s'+1,\dots,s\}$ are the two cyclic
intervals of $\Z/L_0\Z$ determined by $s,s'$.
\end{lemma}

\begin{proof}
$I_1$ and $I_2$ partition $\Z/L_0\Z$. A spatial plaquette has $T(p)=\{r\}$, contained in one of them. A temporal plaquette of layer
$r\notin\{s,s'\}$ has $T(p)=\{r,r+1\}$, and $r,r+1$ lie in the same interval because the only consecutive pairs split between
$I_1$ and $I_2$ are $\{s,s+1\}$ and $\{s',s'+1\}$. If $p,p'\in Y$ are adjacent, the time coordinates of the endpoints of a
common link lie in $T(p)\cap T(p')$, so $p,p'$ lie in the same class. By connectedness all plaquettes of $Y$ lie in the same class.
\end{proof}

\begin{corollary}[Size of connecting components]\label{mg:cor:size}
Let all $L_\lambda\ge3$. Let $S_A,S_B\subseteq E(\Lambda)$ and integers $a_1\le a_2<b_1\le b_2$ with $b_2-a_1<L_0$ be such that the
endpoints of the links of $S_A$ have time coordinates in $\{a_1,\dots,a_2\}$ and those of $S_B$ in $\{b_1,\dots,b_2\}$ (modulo $L_0$). Put
\[
\tau_{\rm in}:=b_1-a_2,\qquad \tau_{\rm out}:=L_0-(b_2-a_1),\qquad \tau:=\min(\tau_{\rm in},\tau_{\rm out})\ \ (\ge1).
\]
Let $X\subseteq\cP$ be such that every link of $E(X)\setminus(S_A\cup S_B)$ lies in at least two plaquettes of $X$, and let $Y$ be a
component of $X$ containing a plaquette of $\cP_{S_A}$ and a plaquette of $\cP_{S_B}$. Then $|Y|\ge g\,\tau$.
\end{corollary}

\begin{proof}
Call the layers $s\in\{a_2,\dots,b_1-1\}$ inner ($\tau_{\rm in}$ of them) and the layers $s'\in\{b_2,\dots,a_1-1+L_0\}$ outer
($\tau_{\rm out}$ of them); as integers, $s<s'<s+L_0$ for every inner $s$ and outer $s'$, so inner and outer layers are distinct.
\emph{These layers are free.} A temporal link of an inner layer $s$ joins $s$ and $s+1$, where $s+1\in\{a_2+1,\dots,b_1\}$ is not in
$\{a_1,\dots,a_2\}$ and $s\in\{a_2,\dots,b_1-1\}$ is not in $\{b_1,\dots,b_2\}$ modulo $L_0$ (because $b_2-a_1<L_0$); so it belongs neither
to $S_A$ nor to $S_B$. The same argument applies to outer layers.
\emph{Times of the plaquettes of $Y$ at the supports.} Let $p\in Y\cap\cP_{S_A}$, containing a link $\ell\in S_A$. If $p$ is spatial,
$T(p)=\{r\}$ with $r\in\{a_1,\dots,a_2\}$; if $p$ is temporal of layer $r$, then $\ell$ is either its temporal link of layer $r$ or a
spatial link at time $r$ or $r+1$, so $r\in\{a_1-1,\dots,a_2\}$. Now let $s$ be inner and $s'$ outer with $Y_s=Y_{s'}=\emptyset$, and
$I_1,I_2$ as in Lemma~\ref{mg:lem:crossing}. Then $T(p)\subseteq I_2=\{s'+1,\dots,s\}$: the times $a_1,\dots,a_2$ lie in $I_2$ because
$s'+1\le a_1+L_0$ and $s\ge a_2$; the time $a_2+1$ occurs in $T(p)$ only if $p$ is temporal of layer $a_2$, which is excluded when
$s=a_2$ (as $Y_s=\emptyset$) and lies in $I_2$ when $s\ge a_2+1$; similarly the time $a_1-1$ occurs only if $p$ is temporal of layer
$a_1-1$, which is excluded when $s'\equiv a_1-1$ and lies in $I_2$ otherwise. In the same way $T(p)\subseteq I_1=\{s+1,\dots,s'\}$ for
every $p\in Y\cap\cP_{S_B}$.
\emph{Conclusion.} If $Y_s=Y_{s'}=\emptyset$ for some inner $s$ and some outer $s'$, Lemma~\ref{mg:lem:crossing} puts all of $Y$ into one
of the disjoint classes $I_1$, $I_2$, which contradicts the previous paragraph, since $Y$ meets both $\cP_{S_A}$ and $\cP_{S_B}$. Hence
$Y_s\neq\emptyset$ for all inner layers, or $Y_{s'}\neq\emptyset$ for all outer layers. In either case Lemma~\ref{mg:lem:layer} (with
$S=S_A\cup S_B$) gives at least $g$ plaquettes of $Y$ in each of these layers, and plaquettes of different layers are different.
\end{proof}

\section{The covariance bound and the mass gap}\label{mg:sec:main}

\begin{theorem}[Covariance bound in finite volume]\label{mg:thm:finite}
Let $d\ge2$, $0<\beta<\beta_0(d)$, and let $\Lambda$ be a periodic lattice with all $L_\lambda\ge3$; let $g$ be the girth of
Lemma~\ref{mg:lem:girth}. Let $A,B$ be bounded measurable functions on $\Omega_\Lambda$ depending only on disjoint link sets $S_A,S_B$
as in Corollary~\ref{mg:cor:size}, with $\tau=\min(\tau_{\rm in},\tau_{\rm out})$. Then, with $q=q(\beta)<1$ (Lemma~\ref{mg:lem:q}) and
$c=c(\beta)$ as in Theorem~\ref{mg:thm:main},
\[
\bigl|\Cov_{\mu_{\Lambda,\beta}}(A,B)\bigr|\le\frac{2\norm A_\infty\norm B_\infty}{1-q}\,c^{|S_A|+|S_B|}\,q^{g\tau}.
\]
\end{theorem}

\begin{proof}
If $S_A$ or $S_B$ is empty, $A$ or $B$ is constant and the covariance vanishes; so let both be non-empty. Let $S:=S_A\cup S_B$ and
$F$ as in Lemma~\ref{mg:lem:cancel}. By Lemmas~\ref{mg:lem:loopexp} and~\ref{mg:lem:cancel},
$2\Cov(A,B)=\sum_X\Phi_F(X)(\Xi(\cP\setminus\overline X)/\Xi(\cP))^2$, where only sets $X\in\cX_S$ contribute in which every link of
$E(X)\setminus S$ lies in at least two plaquettes of $X$ and some component meets both $\cP_{S_A}$ and $\cP_{S_B}$; by
Corollary~\ref{mg:cor:size} these have $|X|\ge n_0:=g\tau$. For each of them: $|\Phi_F(X)|\le4\norm A_\infty\norm B_\infty\,\epsilon(2\beta)^{|X|}$
by Lemmas~\ref{mg:lem:cancel}(i) and~\ref{mg:lem:dup}(a); and by \IIIref{Proposition}{Removing plaquettes}{al:prop:remove} and
Definition~\ref{mg:def:loop},
\[
\Bigl(\frac{\Xi(\cP\setminus\overline X)}{\Xi(\cP)}\Bigr)^2\le e^{2\theta|\overline X|}\le e^{2\theta(K+1)|X|}e^{2\theta J_S}\le e^{2|X|}e^{2\theta J_S},
\]
using $(K+1)\theta<1$ \IIIlab{Lemma}{al:lem:constants}. The number of $X\in\cX_S$ with $|X|=n$ is at most $\binom{J_S+Kn}{n}$: this is
\IIIref{Corollary}{Counting loop polymers}{al:cor:countloop} with $E_C$ replaced by $S$ (its proof uses only $|\cP_S|\le J_S$ and the
adjacency bound $K$; it requires $J_S\ge1$, i.e.\ $S\neq\emptyset$). With $\binom mnt^n\le(1+t)^m$ at $t=1/K$, $\binom{J_S+Kn}n\le(1+\frac1K)^{J_S}(1+\frac1K)^{Kn}K^n\le(1+\frac1K)^{J_S}(eK)^n$.
Altogether
\[
2|\Cov(A,B)|\le4\norm A_\infty\norm B_\infty\,e^{2\theta J_S}\Bigl(1+\frac1K\Bigr)^{J_S}\sum_{n\ge n_0}\bigl(eK\cdot e^2\epsilon(2\beta)\bigr)^n
=\frac{4\norm A_\infty\norm B_\infty\,c^{|S|}}{1-q}\,q^{n_0},
\]
since $e^3K\epsilon(2\beta)=q$ and $(e^{2\theta}(1+\frac1K))^{J_S}=c^{|S|}$.
\end{proof}

\begin{lemma}[Constants]\label{mg:lem:q}
For $d\ge2$ and $0<\beta<\beta_0(d)$: $q(\beta)<\frac{e(2+\epsilon_0)}{K+2}$, which is $<0.909$ for $d=2$ and $<0.389$ for $d\ge3$; and $4\log\frac1{q(\beta)}=4\log\frac1\beta-4\log(4e^3K)+O(\beta)$
as $\beta\downarrow0$. Moreover $\beta_0(d)<\beta_*(d)=\artanh\frac1{6(d-1)}$, the threshold of \cite{II}.
\end{lemma}

\begin{proof}
With $\epsilon=e^{2\beta}-1<\epsilon_0$ (which is the condition $\beta<\beta_0(d)$), $e^{4\beta}-1=\epsilon(2+\epsilon)<\epsilon_0(2+\epsilon_0)$ and
$e^3K\epsilon_0=\frac e{K+2}$; for $d=2$, $K=4$ and $\frac{e(2+\epsilon_0)}6<0.909$; for $d\ge3$, $K\ge12$ and $\frac{e(2+\epsilon_0)}{14}<0.389$. The asymptotics follow from
$e^{4\beta}-1=4\beta+O(\beta^2)$. The last claim is \IIref{Corollary}{Area law in infinite volume}{iv:cor:arealaw} (first line of its proof).
\end{proof}

\begin{proof}[Proof of Theorem~\ref{mg:thm:main}]
The bound on $q$ is Lemma~\ref{mg:lem:q}.
(iii) Since $\beta<\beta_*(d)$, the infinite-volume Gibbs measure $\mu_\beta$ of \IIref{Theorem}{Uniqueness and convergence}{iv:thm:unique}
exists and is the limit of the periodised measures $\tilde\mu_{\Lambda_k,\beta}$ on every bounded measurable local function, for
every sequence of periodic lattices with all sides tending to infinity. Take such $\Lambda_k$ with spatial sides $\ge4$. For $k$ large,
the reduction map is injective on the links of $S_A\cup S_B$, the functions $A\circ\iota$, $B\circ\iota$ on $\Omega_{\Lambda_k}$ (in the notation
of \cite{II}) depend on the images of $S_A$, $S_B$, which satisfy the hypotheses of Theorem~\ref{mg:thm:finite} with the same $a_i,b_i$,
$g=4$, $\tau_{\rm in}=b_1-a_2$ and $\tau_{\rm out}\to\infty$. Theorem~\ref{mg:thm:finite} bounds
$\Cov_{\tilde\mu_{\Lambda_k,\beta}}(A,B)$ uniformly in $k$, and the left side converges to $\Cov_{\mu_\beta}(A,B)$.

(ii) Let $F\in\cA_+$ in the notation of \IIref{Definition}{Reflection, shifts, half-space algebras}{iv:def:reflection}, depending only on a
finite set of links whose endpoints have time coordinates in $[0,n_F]$. Then $\Theta F$ depends only on links with endpoint times in
$[-n_F,0]$ and $S_tF$ on links with endpoint times in $[t,t+n_F]$, with the same number of links. For $t\ge1$ the two sets are disjoint
and (iii) gives, by $\theta$- and translation invariance of $\mu_\beta$ (so that $\mu_\beta(\Theta F)\mu_\beta(S_tF)=|\mu_\beta(F)|^2$),
\[
\bigl|\mu_\beta(\Theta F\cdot S_tF)-|\mu_\beta(F)|^2\bigr|=\bigl|\Cov_{\mu_\beta}(\Theta F,S_tF)\bigr|\le C_F\,q^{4t}\qquad(t\ge1),
\]
with $C_F$ independent of $t$ (here $b_1-a_2=t$); for $t=0$ the left side is at most $2\norm F_\infty^2$. So the hypothesis of
\IIlab{Theorem}{iv:thm:OS}(iii) holds with $m=4\log(1/q)$ and the constant $\max(C_F,2\norm F_\infty^2)$, and
$m(\mu_\beta)\ge4\log(1/q)$.

(i) Let $A$ be a bounded gauge-invariant function on $\Omega_s=G^{E(\Lambda_s)}$ and $\Lambda=(\Z/L_0\Z)\times\Lambda_s$ with $L_0\ge3$.
The functions $\overline A(t)$ and $A(0)$ depend on the spatial links of the slices $t$ and $0$: in Corollary~\ref{mg:cor:size} take
$a_1=a_2=0$, $b_1=b_2=t$; for $1\le t\le L_0/2$, $\tau=\min(t,L_0-t)=t$, and $g=4$ since all $L_i\ge4$. Theorem~\ref{mg:thm:finite} gives
$|\langle\overline A(t)A(0)\rangle-\langle\overline A(t)\rangle\langle A(0)\rangle|\le C_Aq^{4t}$ with $C_A$ depending on $\norm A_\infty$ and
$|E(\Lambda_s)|$ only, not on $L_0$; for $t=0$ the left side is at most $2\norm A_\infty^2$. \Flab{Corollary}{cor:gapfromcorr} (whose proof
lets $L_0\to\infty$ at fixed $t$, so that the bound is needed only for large $L_0$) gives $m_{\Lambda_s}(\beta)\ge4\log(1/q)$.
\end{proof}

\section{Two dimensions; discussion}\label{mg:sec:discussion}

\begin{proposition}[Two dimensions]\label{mg:prop:2d}
Let $d=2$ and $0<\beta<\beta_0(2)$, so that $K=4$ and $q(\beta)<0.909$ (Lemma~\ref{mg:lem:q}).
\begin{enumerate}[label=(\roman*)]
\item For every $L_1\ge3$, the spatial lattice $\Lambda_s=\Z/L_1\Z$ has $m_{\Lambda_s}(\beta)\ge L_1\log(1/q(\beta))$.
\item In infinite volume, $\Cov_{\mu_\beta}(A,B)=0$ for all $A,B$ as in Theorem~\ref{mg:thm:main}(iii) (time ranges with $a_2<b_1$);
hence $m(\mu_\beta)=\infty$.
\end{enumerate}
\end{proposition}

\begin{proof}
By Lemma~\ref{mg:lem:girth}, $g=L_1$, and Theorem~\ref{mg:thm:finite} gives the bound with $q^{L_1\tau}$. (i) follows as in the proof of
Theorem~\ref{mg:thm:main}(i), with $L_1$ in place of $4$. For (ii), the bound for the periodised observables on tori with $L_0,L_1\to\infty$
tends to $0$, since $\tau\ge1$ and $g=L_1\to\infty$, and the torus covariances converge to $\Cov_{\mu_\beta}(A,B)$ as in the proof of
Theorem~\ref{mg:thm:main}(iii) (note $\beta_0(2)<\beta_*(2)$, Lemma~\ref{mg:lem:q}); then \IIlab{Theorem}{iv:thm:OS}(iii) applies with every
$m>0$.
\end{proof}

\begin{remark}[Consistency with the exact solution]
In $d=2$ the transfer matrix on a spatial circle of length $L_1$ is diagonal in the characters of the holonomy, with eigenvalue ratios
$(c_r/(d_rc_0))^{L_1}$ (\texttt{lattice/ymlattice/transfer2d.py}, tested against the Monte Carlo); so the exact finite-volume gap is
$L_1\log(1/u(\beta))$ for $SU(2)$, linear in $L_1$, and the infinite-volume theory has no propagating states. Proposition~\ref{mg:prop:2d}
reproduces both features (with $q$ in place of $u$). It is also consistent with \IIref{Proposition}{Gauge invariance of the ground-state
chain}{iv:prop:phys}: in $d=2$ the spatial links of a time slice form a line, gauge-invariant functions of finitely many of them are constant,
and $P$ acts trivially on the physical sector. The difference between $d=2$ and $d\ge3$ is the girth of the spatial graph ($L_1$ versus
$4$): in $d\ge3$ a closed surface can turn a corner in every time layer with four plaquettes, the minimal glueball.
\end{remark}

\begin{remark}[The strong-coupling test, sharpness, and the $U(1)$ test]\label{mg:rem:tests}
(i) Principle~3 of the research program requires the correct asymptotics $m(\beta)\sim4\log(1/u(\beta))$. Theorem~\ref{mg:thm:main}
gives the lower bound $m\ge4\log(1/\beta)-O(1)$; this is the lower half. For the upper half: since $P\ge0$
\IIref{Theorem}{Transfer operator, ground-state chain and mass gap}{iv:thm:gap}(ii), the effective masses of any correlator decrease to a
limit $\ge m(\mu_\beta)$ \IIref{Remark}{Consequences for the numerics}{iv:rem:numerics}, so a computation of a single correlator (e.g.\ of two
plaquettes) at two fixed separations, to leading order with a remainder bound, would give $m\le4\log(1/u(\beta))+O(\beta)$. This is
expected but not proved here. (ii) The constant $-4\log(4e^3K)$ reflects the crude counting of \cite{III} and the crude bound
$\norm{g_p}_\infty\le e^{4\beta}-1$: for $SU(N)$ and $U(1)$, $c_0=1+O(\beta^2)$ gives $\norm{g_p}_\infty\le2\beta+O(\beta^2)$, which improves the
constant by $4\log2$, and counting only surfaces without free edges, as in \IIIref{Remark}{Constants, sharpness,
improvements}{al:rem:constants}(ii), would improve it further. (iii) The hypothesis $L_i\ge4$ in (i) is needed: if a spatial side equals $3$,
Theorem~\ref{mg:thm:finite} holds with $g=3$, and closed strips of three plaquettes per layer around that side (torelons, carrying the flux of
a Polyakov loop) are expected (heuristic) to give states of energy $\approx3\log(1/u(\beta))$, below $4\log(1/q(\beta))$ for very small
$\beta$. (iv) Nothing here is non-abelian: the argument applies verbatim to $G=U(1)$, whose $d=4$ theory is expected to be massless at large
$\beta$ (proved in \cite{FS82} for the Villain action and asserted there for Wilson's; \IIIref{Remark}{The $U(1)$ test}{al:rem:u1test}). The
result therefore concerns the strong-coupling endpoint only; it shows that the endpoint is controlled with the right leading behaviour of the
lower bound, which is what a renormalisation-group argument would need to reach.
\end{remark}

\subsection*{Summary of results}

\begin{center}
\small
\begin{tabular}{@{}p{0.6\textwidth}p{0.2\textwidth}l@{}}
\toprule
Statement & Where & Status\\
\midrule
Duplicated system; loop expansion with $\Xi^2$; cancellation of non-connecting terms & Lem.~\ref{mg:lem:dup}, \ref{mg:lem:loopexp}, \ref{mg:lem:cancel} & Theorem\\
Surfaces without free edges have $\ge g$ temporal plaquettes per crossed layer ($g=4$ for $d\ge3$, spatial sides $\ge4$) & Lem.~\ref{mg:lem:girth}, \ref{mg:lem:layer}, \ref{mg:lem:crossing}, Cor.~\ref{mg:cor:size} & Theorem\\
Covariance bound $q^{g\tau}$ in finite and infinite volume & Thm.~\ref{mg:thm:finite}, \ref{mg:thm:main}(iii) & Theorem\\
$m(\mu_\beta)\ge4\log(1/q(\beta))=4\log(1/\beta)-O(1)$, and $m_{\Lambda_s}(\beta)\ge4\log(1/q(\beta))$ for spatial sides $\ge4$; $d\ge3$, $\beta<\beta_0(d)$ & Thm.~\ref{mg:thm:main} & Theorem\\
$d=2$: gap $\ge L_1\log(1/q)$; infinite gap in infinite volume & Prop.~\ref{mg:prop:2d} & Theorem\\
Matching upper bound $m\le4\log(1/u(\beta))+O(1)$ & Rem.~\ref{mg:rem:tests} & open (expected)\\
\bottomrule
\end{tabular}
\end{center}

\begin{thebibliography}{OBS85}

\bibitem[F]{F}
\emph{Lattice Yang--Mills theory: foundations for the mass gap problem}, Part~I of the project ``Solving Yang--Mills and the Mass
Gap'', \texttt{theory/lattice-foundations.tex} (this repository), living document.

\bibitem[II]{II}
\emph{Lattice Yang--Mills theory at strong coupling II: the infinite-volume theory, its transfer operator and its mass gap},
\texttt{theory/infinite-volume-strong-coupling.tex} (this repository), living document, version of 9~November~2026.

\bibitem[III]{III}
\emph{Lattice Yang--Mills theory at strong coupling III: confinement (area law)}, \texttt{theory/area-law-strong-coupling.tex}
(this repository), living document, revised 9~November~2026.

\bibitem[FS82]{FS82}
J.~Fröhlich, T.~Spencer, Massless phases and symmetry restoration in abelian gauge theories and spin systems,
Commun.\ Math.\ Phys.\ 83 (1982), 411--454. (verify)

\bibitem[OB85]{OB85}
M.~O'Carroll, W.~D.~Barbosa, J.\ Math.\ Phys.\ 26 (1985): glueball masses in $3+1$ dimensional lattice gauge theory at strong
coupling (title and content to be checked; only an abstract has been seen). (verify)

\bibitem[OBS85]{OBS85}
M.~O'Carroll, G.~Braga, R.~Schor, Commun.\ Math.\ Phys.\ 97 (1985): glueball masses in $2+1$ dimensional lattice gauge theory at
strong coupling (title and content to be checked; only an abstract has been seen). (verify)

\bibitem[Sch83]{Sch83}
R.~Schor, Existence of glueballs in strongly coupled lattice gauge theories, Nucl.\ Phys.\ B222 (1983), 71--82. (verify)

\end{thebibliography}

\end{document}
